%=====================================================================
\chapter{Sixteen-Week Teaching Schedule}\label{app:schedule}
%=====================================================================
\normalfigures

A mapping from the 32 chapters onto a standard 16-week semester of two lecture
hours plus one lab hour per week, with assessment points. Adjust to your own
calendar; the dependencies are what matter.

\section{The Schedule}

\rowcolors{2}{white}{lightbg}
\begin{longtable}{@{}p{1.1cm}p{3.4cm}p{3.4cm}p{3.1cm}p{2.6cm}@{}}
\toprule
\rowcolor{primary!15}
\textbf{Wk} & \textbf{Lecture topic} & \textbf{Chapters} & \textbf{Lab} & \textbf{Assessment} \\
\midrule
\endhead
1 & The landscape and the lifecycle & 1, 2 & Meet a dataset; write a frame & --- \\
2 & Data and the mathematical toolkit & 3, 4 & Ten-minute audit; distance \& descent & --- \\
3 & Describing data; probability & 5, 6 & Centre, spread, shape; base-rate simulation & Quiz 1 (Ch 1--4) \\
4 & Wrangling and cleaning & 7 & Leak-proof pipeline & --- \\
5 & Exploratory analysis & 8 & Standard EDA pass & \textbf{Assignment 1 set} \\
6 & Inference and causality & 9, 10 & Intervals, tests, confounder simulation & --- \\
7 & ML fundamentals; regression & 11, 12 & Split, baseline, diagnose; fit \& regularise & \textbf{Assignment 1 due} \\
8 & Classification & 13 & Five classifiers compared & \textbf{Midterm (Ch 1--13)} \\
9 & \textbf{Evaluation} & 14 & Evaluate honestly; cost-based threshold & --- \\
10 & Features and ensembles; unsupervised & 15, 16 & Engineer, ensemble, interpret; cluster \& detect & \textbf{Assignment 2 set} \\
11 & Neural networks and architectures & 17, 18 & Network from scratch; transfer learning & --- \\
12 & Time series; text and NLP & 19, 20 & Decompose \& walk-forward; TF-IDF classifier & \textbf{Assignment 2 due} \\
13 & Generative and agentic AI & 21, 22 & Minimal RAG; agent loop with guardrails & Quiz 2 (Ch 14--22) \\
14 & Data engineering; cloud and edge & 23, 24 & Data contracts; containerise a model & \textbf{Capstone set} \\
15 & \textbf{Data science for IoT and cybersecurity} & 25, 26 & Telemetry to alert; per-entity baselines \& evasion test & --- \\
16 & MLOps; ethics; practice \& capstone & 27, 28, 29--32 & Monitored service; fairness measurement & \textbf{Capstone due} \\
\bottomrule
\end{longtable}

\begin{alertbox}[If You Must Cut Something]
Cut \emph{depth}, never \emph{Chapter 14}. Evaluation is the hinge of the whole
course: a student who cannot evaluate honestly will misreport every subsequent
result. If time is short, compress Chapters 18 and 20--22 --- they are the most
self-contained --- and protect Chapters 2, 7, 11, 14 and 28, which everything else
depends on.
\end{alertbox}

\section{Dependency Map}

\dsfigh{%
\begin{tikzpicture}[font=\scriptsize,
  n/.style={rounded corners=2pt, draw=#1, fill=#1!10, text=#1!50!black,
            minimum width=1.15cm, minimum height=0.5cm, font=\tiny\bfseries},
  crit/.style={rounded corners=2pt, draw=accent, line width=1.4pt, fill=accent!16,
               text=accent!70!black, minimum width=1.15cm, minimum height=0.5cm,
               font=\tiny\bfseries},
  a/.style={-{Stealth[length=1.6mm]}, gray!55, line width=0.6pt}]

% Part I row
\foreach \i/\c in {1/1,2/2,3/3,4/4,5/5,6/6}{
  \pgfmathsetmacro{\xx}{(\i-1)*1.45}
  \ifnum\c=2
    \node[crit] (c\c) at (\xx,2.4) {\c};
  \else
    \node[n=primary] (c\c) at (\xx,2.4) {\c};
  \fi}
\node[font=\tiny, anchor=west, text=primary] at (8.9,2.4) {Part I --- foundations};

% Part II
\foreach \i/\c in {1/7,2/8,3/9,4/10}{
  \pgfmathsetmacro{\xx}{(\i-1)*1.45}
  \ifnum\c=7
    \node[crit] (c\c) at (\xx,1.55) {\c};
  \else
    \node[n=secondary] (c\c) at (\xx,1.55) {\c};
  \fi}
\node[font=\tiny, anchor=west, text=secondary] at (8.9,1.55) {Part II --- data to insight};

% Part III
\foreach \i/\c in {1/11,2/12,3/13,4/14,5/15,6/16}{
  \pgfmathsetmacro{\xx}{(\i-1)*1.45}
  \ifnum\c=14
    \node[crit] (c\c) at (\xx,0.7) {\c};
  \else
    \ifnum\c=11
      \node[crit] (c\c) at (\xx,0.7) {\c};
    \else
      \node[n=greenhl] (c\c) at (\xx,0.7) {\c};
    \fi
  \fi}
\node[font=\tiny, anchor=west, text=greenhl!50!black] at (8.9,0.7) {Part III --- ML core};

% Part IV
\foreach \i/\c in {1/17,2/18,3/19,4/20,5/21,6/22}{
  \pgfmathsetmacro{\xx}{(\i-1)*1.45}
  \node[n=purplehl] (c\c) at (\xx,-0.15) {\c};}
\node[font=\tiny, anchor=west, text=purplehl] at (8.9,-0.15) {Part IV --- deep learning};

% Part V
\foreach \i/\c in {1/23,2/24,3/25,4/26,5/27}{
  \pgfmathsetmacro{\xx}{(\i-1)*1.45}
  \node[n=tealhl] (c\c) at (\xx,-1.0) {\c};}
\node[font=\tiny, anchor=west, text=tealhl] at (8.9,-1.0) {Part V --- systems};

% Part VI
\foreach \i/\c in {1/28,2/29,3/30,4/31,5/32}{
  \pgfmathsetmacro{\xx}{(\i-1)*1.45}
  \ifnum\c=28
    \node[crit] (c\c) at (\xx,-1.85) {\c};
  \else
    \node[n=goldhl] (c\c) at (\xx,-1.85) {\c};
  \fi}
\node[font=\tiny, anchor=west, text=goldhl!50!black] at (8.9,-1.85) {Part VI --- practice};

% key dependencies
\draw[a] (c2) -- (c7);
\draw[a] (c6) -- (c9);
\draw[a] (c4) -- (c11);
\draw[a] (c14) -- (c25);
\draw[a] (c14) -- (c26);
\draw[a] (c16) -- (c26);
\draw[a] (c19) -- (c25);
\draw[a] (c24) -- (c25);
\draw[a] (c14) -- (c28);

\node[anchor=west, align=left, text width=6.0cm, font=\tiny, text=accent!75!black]
     at (8.9,-3.0)
     {\textbf{Outlined chapters are load-bearing.} 2 (framing), 7 (leakage),
      11 (generalisation), 14 (evaluation) and 28 (ethics) are relied on by
      everything downstream. Arrows show the strongest cross-part dependencies ---
      note that Chapters 25 and 26 each pull from three earlier parts, which is why
      they sit where they do.};
\end{tikzpicture}}%
{Chapter dependency map. The outlined chapters cannot be skipped; the arrows show why
Chapters 25 and 26 are placed in Part V rather than earlier.}{dependency-map}

\section{Assessment Design}

\rowcolors{2}{white}{defbg}
\begin{longtable}{@{}p{3.2cm}p{1.6cm}p{3.0cm}p{5.8cm}@{}}
\toprule
\rowcolor{primary!15}
\textbf{Component} & \textbf{Weight} & \textbf{Covers} & \textbf{Assesses} \\
\midrule
\endhead
Quiz 1 (week 3) & 5\% & Ch 1--4 & Vocabulary and the analytics levels; a low-stakes early check \\
Assignment 1 (wk 5--7) & 15\% & Ch 3, 5, 7, 8 & An audit and EDA on a messy dataset, with documented cleaning decisions \\
Midterm (week 8) & 20\% & Ch 1--13 & Conceptual understanding; hand calculations; framing a stated problem \\
Assignment 2 (wk 10--12) & 20\% & Ch 11--16 & A complete supervised project with baseline, correct validation and cost-based threshold \\
Quiz 2 (week 13) & 5\% & Ch 14--22 & Evaluation, architectures, generative and agentic concepts \\
Capstone (wk 14--16) & 35\% & All & The four deliverables of \chref{capstone}, marked on that chapter's rubric \\
\bottomrule
\end{longtable}

\begin{conceptbox}[Two Assessment Design Notes for the Instructor]
\textbf{Set Assignment 2 on an imbalanced dataset.} Students who report accuracy on
it have not understood \chref{evaluation}, and a single well-chosen dataset
discriminates better than any exam question.

\textbf{Award marks in the capstone for a null result honestly reported.} A student
who establishes that the signal is not present, documents why, and recommends
stopping has demonstrated the judgement of \chref{projectmgmt} --- and rewarding
that is the only way to stop students overclaiming.
\end{conceptbox}

\section{Lab Materials}

Every ``Try It Yourself'' box in the handout is a lab exercise. The runnable
extracts live in \texttt{code/} in this folder. Labs assume Python as primary and R
where the department's existing Quarto materials already cover the topic
(Appendix~B).

\rowcolors{2}{white}{lightbg}
\begin{longtable}{@{}p{2.0cm}p{5.6cm}p{6.4cm}@{}}
\toprule
\rowcolor{primary!15}
\textbf{Week} & \textbf{Lab} & \textbf{Deliverable} \\
\midrule
\endhead
1--2 & Meet a dataset; write a frame; ten-minute audit & A completed six-question frame and an audit report \\
3 & Centre, spread, shape; base-rate simulation & A short note on why the detector's alerts are mostly wrong \\
4--5 & Leak-proof pipeline; standard EDA pass & A cleaned dataset plus five charts with one-sentence findings \\
6--7 & Intervals and tests; confounder simulation; fit and regularise & A regression with residual diagnostics \\
8--9 & Five classifiers; evaluate honestly & A PR curve and a cost-justified threshold \\
10 & Engineer, ensemble, interpret; cluster and detect & Permutation importances and a validated clustering \\
11--12 & Network from scratch; transfer learning; decompose and walk-forward & A learning-curve diagnosis and a temporal model beating naive \\
13 & Minimal RAG; agent loop with guardrails & A grounded assistant with a documented evaluation set \\
14--15 & Data contracts; containerise; telemetry to alert; per-entity baselines & A running container and an evasion test on your own model \\
16 & Monitored service; fairness measurement & A model card and a four-layer monitoring plan \\
\bottomrule
\end{longtable}
